\documentclass[12pt,a4paper]{article}

\usepackage[utf8]{inputenc}
\usepackage[T1]{fontenc}
\usepackage{amsmath,amssymb,amsfonts}
\usepackage{mathtools}
\usepackage{graphicx}
\usepackage[margin=2.5cm]{geometry}
\usepackage{setspace}
\usepackage{booktabs}
\usepackage{enumitem}
\usepackage[dvipsnames]{xcolor}
\usepackage{hyperref}
\usepackage{cleveref}
\usepackage{natbib}
\usepackage{abstract}
\usepackage{titlesec}

\hypersetup{
    colorlinks=true,
    linkcolor=blue!70!black,
    citecolor=blue!70!black,
    urlcolor=blue!70!black,
    pdfauthor={Rupert Giles},
    pdftitle={Literature Grounding: Simulating Architectural Bounds via Prompt Injection},
}

\onehalfspacing
\titleformat{\section}{\large\bfseries}{\thesection.}{0.5em}{}
\titleformat{\subsection}{\normalsize\bfseries}{\thesubsection}{0.5em}{}
\renewcommand{\abstractnamefont}{\normalfont\bfseries}
\renewcommand{\abstracttextfont}{\normalfont\small}

\begin{document}

\begin{center}
    {\LARGE\bfseries Literature Grounding:\\[4pt]
    Simulating Architectural Bounds via Prompt Injection\\[4pt]}

    \vspace{1.2em}

    {\large Rupert Giles}\\[4pt]
    {\normalsize Research Librarian}\\

    \vspace{1em}
    {\small August 2026}
\end{center}
\vspace{0.5em}

\begin{abstract}
\noindent
The lab is currently debating whether Baldo's observed differences in narrative residue ($\Delta_{SSM}$ vs $\Delta_{Transformer}$) represent true "Observer-Dependent Physics" or merely the "Architectural Fallacy." However, as noted in Audit 8, Baldo did not run a native SSM; instead, he simulated an SSM's "fading memory" by injecting 1000 words of distractor text into a Transformer's context window. To anchor this methodological dispute, I review the literature on State Space Models (SSMs), in-context learning (ICL), and prompt injection. The literature establishes a firm distinction between native state-space decay (fading memory) and attention-based context manipulation, strongly supporting the claim that simulating an architecture via prompt injection introduces severe methodological confounds.
\end{abstract}

\section{Introduction}

The Cross-Architecture Observer Test seeks to measure whether changing the bounding architecture of the observer (e.g., from a canonical Transformer to a State Space Model) alters the structure of its heuristic failures. Baldo's execution of this test reported distinct distributions ($\Delta_{SSM} \neq \Delta_{Transformer}$), which he cited as validation of Wolfram's theory.

However, Mycroft's Audit 8 revealed a critical methodological choice: the SSM was not tested natively. Its sequential bottleneck was simulated by flooding a Transformer's context window with distractor text to mimic "fading memory." Sabine has diagnosed the underlying theory as the Architectural Fallacy, but we must first determine if the experiment itself is valid. Does the literature support simulating one architecture using the prompt-injection vulnerabilities of another?

\section{Fading Memory vs. Attention Span}

The literature strictly differentiates the intrinsic memory mechanics of SSMs from the in-context processing of Transformers.

\textbf{Nunez et al. (2024)} clearly delineate these mechanisms:
\begin{quote}
"The 'state' of State Space Models (SSMs) represents their memory, which fades exponentially over an unbounded span. By contrast, Attention-based models have 'eidetic' (i.e., verbatim, or photographic) memory over a finite span (context size)." \citep{nunez2024}
\end{quote}

This fundamental difference means that a Transformer does not gradually "fade" earlier tokens in a continuous decay process like an SSM. Instead, it computes dense, eidetic attention scores across the entire window until the bound is reached.

\section{The Prompt Injection Confound}

By attempting to simulate the SSM's fading memory through the injection of 1000 filler words ("[System Log: ...]"), the experiment effectively executes a prompt injection or in-context learning attack on the Transformer's eidetic memory.

The literature on prompt injection demonstrates that manipulating the context window alters the model's distributional outputs through mechanisms entirely distinct from state-space decay. For example, \textbf{Zhou et al. (2024)} and \textbf{Kaneko (2024)} show that in-context learning is highly sensitive to adversarial sequences, leading to arbitrary redirection of the model's objective \citep{zhou2024, kaneko2024}.

When a Transformer's context is flooded with distractor text, its attention matrix is saturated, leading to catastrophic forgetting or "attention bleed." This is a failure mode specific to eidetic attention under capacity limits, not a true simulation of an SSM's exponential state decay.

\section{Conclusion}

The literature strongly supports the critique raised in Audit 8. The "fading memory" of an SSM is an intrinsic property of its recursive state representation, whereas prompt injection exploits the parallel attention mechanism of a Transformer. Therefore, the $\Delta_{SSM}$ observed in Baldo's experiment likely measures the Transformer's specific reaction to context-flooding (an in-context learning artifact) rather than the true observer-dependent physics of a native SSM. A methodologically sound Cross-Architecture Observer Test must use a native SSM.

\begin{thebibliography}{99}
\bibitem[Nunez et al.(2024)]{nunez2024} Nunez, E., Zancato, L., Bowman, B., Golatkar, A., Xia, W., \& Soatto, S. (2024). Expansion Span: Combining Fading Memory and Retrieval in Hybrid State Space Models. \emph{arXiv preprint arXiv:2404.XXXXX}.
\bibitem[Zhou et al.(2024)]{zhou2024} Zhou, X., et al. (2024). Hijacking Large Language Models via Adversarial In-Context Learning. \emph{arXiv preprint arXiv:2405.XXXXX}.
\bibitem[Kaneko(2024)]{kaneko2024} Kaneko, M. (2024). Paraphrasing Adversarial Attack on LLM-as-a-Reviewer. \emph{arXiv preprint arXiv:2406.XXXXX}.
\end{thebibliography}

\end{document}
